\section{Определения}

Комплексные числа. Алгебраическая и тригонометрическая формы.

Ряды с комплексными членами. Предел и непрерывность функции комплексной переменной. Сфера Римана.

\subsection{Ряды}

\begin{definition}
Говорят, что функциональная последовательность $f_k : X \to \mathbb{C}$ сходится \textbf{поточечно} (или просто сходится) на множестве $D \subset X$, если
\[
\forall x \in D \quad \exists \lim_{k \to \infty} f_k(x) \in \mathbb{C}.
\]
Множество $D$ при этом называется \textbf{множеством (поточечной) сходимости} функциональной последовательности $f_k$.
\end{definition}

\begin{definition}
Говорят, что функциональный ряд с общим членом $f_k : X \to \mathbb{C}$ сходится \textbf{поточечно} (или просто сходится) на множестве $D \subset X$, если
\[
\forall x \in D \quad \sum_{k=1}^{\infty} f_k(x) \text{ сходится.}
\]
Множество $D$ при этом называется \textbf{множеством (поточечной) сходимости} функционального ряда с общим членом $f_k$.
\end{definition}

\begin{definition}
Говорят, что последовательность $f_k : X \to \mathbb{C}$ сходится к функции $f$ на множестве $D \subset X$ \textbf{равномерно}, если
\[
\forall \varepsilon > 0 \; \exists k_0 \in \mathbb{N} : \; \forall k > k_0 \; \forall x \in D \quad |f_k(x) - f(x)| < \varepsilon.
\]
Обозначают это так:
\[
f_k \xRightarrow[k \to \infty]{D} f.
\]
\end{definition}

\begin{definition}
Говорят, что функциональный ряд с общим членом $f_k : X \to \mathbb{C}$ сходится \textbf{равномерно} на множестве $D \subset X$, если последовательность его частичных сумм сходится равномерно на $D$.
\end{definition}

\subsection{Дифференцируемость}

Рассмотрим функцию $f: \mathbb{C}\supset D \to \mathbb{C}$ на комплексной плоскости как отображение $\mathbb{R}^2\supset D \to \mathbb{R}^2$, сопоставляющее каждой точке $z = x + iy$ точку
\[
f(z) = f(x, y) = u(x, y) + i v(x, y) = \operatorname{Re} f(x, y) + i \operatorname{Im} f(x, y).
\]

\begin{definition}
Функция $f(x, y) = u(x, y) + i v(x, y)$, определённая в окрестности точки $z_0 = x_0 + i y_0$, называется \textbf{$\mathbb{R}$-дифференцируемой} в точке $z_0$, если функции $u(x, y)$ и $v(x, y)$ дифференцируемы в точке $(x_0, y_0)$ как функции от двух вещественных переменных $x$ и $y$.
\end{definition}

\begin{definition}
Функция $f$, определённая в окрестности точки $z_0$, называется \textbf{$\mathbb{C}$-дифференцируемой (голоморфной)} в точке $z_0$, если найдется комплексное число $a$ такое, что в окрестности точки $z_0$ имеет место представление
\[
\Delta f = f(z) - f(z_0) = a \cdot \Delta z + o(\Delta z),
\]
где $\Delta z = z - z_0$, $\Delta z\to 0$.
\end{definition}

\begin{definition}
Функция, голоморфная во всей комплексной плоскости $\mathbb{C}$, называется \textbf{целой}.
\end{definition}

\begin{definition}
Функция $f: D \to \mathbb{C}$, определённая в области $D \subset \mathbb{C}$, называется \textbf{аналитичной} в точке $a \in D$, если существует степенной ряд
\[
\sum_{n=0}^{\infty} c_n (z - a)^n,
\]
сходящийся в некоторой окрестности точки $a$, сумма которого равна $f(z)$ в этой окрестности.

Если $f$ аналитична в каждой точке области $D$, то говорят, что $f$ аналитична в $D$.
\end{definition}

\begin{definition}
Пусть функция $f$ $\mathbb{R}$-дифференцируема в точке $z_0$. Отображение окрестности этой точки в $\mathbb{C}$, задаваемое функцией $f$, называется \textbf{конформным в точке} $z_0$, если его дифференциал $df(z_0)$, рассматриваемый как линейное отображение плоскости $\mathbb{R}^2$ на себя, невырожден (т.е. взаимно однозначен) и является композицией поворота и растяжения.
\end{definition}

Способы вычисления интеграла при заданной параметризации.

Пути?

Длина кривой

\subsection{Криволинейный интеграл}

\begin{definition}
Пусть $\gamma: I \to \mathbb{C}$ --- кусочно-гладкий путь, параметризованный отрезком $I = $. Рассмотрим функцию $f: \gamma(I) \to \mathbb{C}$ такую, что композиция $f \circ \gamma$ непрерывна на $I$. Тогда комплекснозначная функция
\[
g(t) = f(\gamma(t)) \dot{\gamma}(t)
\]
интегрируема по Риману на $I = [\alpha, \beta]$, и её интеграл
\[
\int_{\gamma} f(z) \, dz := \int_{\alpha}^{\beta} f(\gamma(t)) \dot{\gamma}(t) \, dt
\]
называется \textbf{интегралом от функции $f$ вдоль пути $\gamma$}.
\end{definition}

\begin{definition}
\textbf{Первообразной} функции $f$ в области $D \subset \mathbb{C}$ называется голоморфная в $D$ функция $F$ такая, что
\[
F'(z) = f(z) \quad \text{для всех } z \in D.
\]
\end{definition}

\subsection{Нули функции}

\begin{definition}
Пусть функция $f$ голоморфна в области $D \subset \mathbb{C}$. Точка $a \in D$ называется \textbf{изолированным нулём} функции $f$, если:
\begin{enumerate}
    \item $f(a) = 0$,
    \item Существует проколотая окрестность $\mathring{U}(a)$, в которой $f(z) \neq 0$.
\end{enumerate}
\end{definition}

\begin{definition}
Пусть функция $f$ голоморфна в окрестности точки $a \in \mathbb{C}$, $f(a) = 0$, но $f$ не равна тождественно нулю в любой проколотой окрестности $a$. Число
\[
n = \min\{m \geq 1 : c_m \neq 0\} = \min\{m \geq 1 : f^{(m)}(a) \neq 0\}
\]
называется \textbf{порядком нуля} голоморфной функции $f$ в точке $a$.

Иначе говоря, порядок нуля $f$ в точке $a$ --- это номер $n$ первой отличной от нуля производной $f^{(n)}(a)$, или, эквивалентно, это единственное натуральное число $n$ такое, что в некоторой окрестности $U$ точки $a$ справедливо представление
\[
f(z) = (z - a)^n g(z)
\]
для некоторой $g \in \mathcal{O}(U)$, $g(a) \neq 0$.
\end{definition}

\subsection{Особые точки}

\begin{definition}
Точка $a \in \mathbb{C}$ называется \textbf{изолированной особой точкой} для функции $f(z)$, если $f$ голоморфна в некоторой проколотой окрестности $V = \{0 < |z - a| < \varepsilon\}, \varepsilon > 0$, точки $a$.

Изолированная особая точка $a$ функции $f$ называется:
\begin{enumerate}
    \item \textbf{устранимой}, если существует (конечный) предел
    \[
    \lim_{z \to a} f(z) \in \mathbb{C};
    \]
    \item \textbf{полюсом}, если существует
    \[
    \lim_{z \to a} f(z) = \infty;
    \]
    \item \textbf{существенно особой точкой} во всех остальных ситуациях, т.е. когда не существует (ни конечного, ни бесконечного) предела $f(z)$ при $z \to a$.
\end{enumerate}
\end{definition}

\begin{definition}
Точка $a = \infty$ называется \text{изолированной особой точкой} для функции $f$, если $f \in \mathcal{O}(\{|z| > R\})$ для некоторого $R > 0$.
\end{definition}

\subsection{Вычет}

Пусть функция $f$ голоморфна в проколотой окрестности $V = \{0 < |z - a| < \varepsilon\}$ точки $a \in \mathbb{C}$, так что $a$ является её изолированной особенностью.

\begin{definition}
\textbf{Вычетом} функции $f$ в изолированной особой точке $a \in \mathbb{C}$ называется число
\[
\operatorname{res}_a f = \frac{1}{2\pi i} \int_{|\zeta - a| = r} f(\zeta) \, d\zeta, \quad \text{где } 0 < r < \varepsilon
\]
(по теореме Коши этот интеграл не зависит от выбора $r$).
\end{definition}

\begin{definition}
\textbf{Вычетом} функции $f$ в бесконечности называется число
\[
\operatorname{res}_{\infty} f = \frac{1}{2\pi i} \int_{\gamma_R^{-}} f(z) \, dz,
\]
где интеграл берётся по окружности $\gamma_R = \{ |z| = R \}$ достаточно большого радиуса $R > R_0$, ориентированной \emph{по часовой стрелке} (отрицательное направление).
\end{definition}

\begin{definition}
\textbf{Логарифмическим вычетом} функции $f$ в точке $a$ называется вычет её логарифмической производной
\[
\frac{f'(z)}{f(z)}
\]
в этой точке. Иными словами, логарифмический вычет $f$ в точке $a$ равен
\[
\operatorname{res}_a \frac{f'(z)}{f(z)} = \frac{1}{2\pi i} \int_{|z-a|=\rho} \frac{f'(z)}{f(z)} \, dz
\]
для любого достаточно малого $\rho > 0$.
\end{definition}

\section{Теоремы}

Геометрический смысл операций над комплексными числами, корень.

Свойства рядов с комплексными членами.

Геометрический смысл модуля и аргумента производной функции f (z) = az, a ∈ C.

Элементарные функции и их свойства. Геометрия экспоненты.

Интеграл γ z n dz при целых n.

Формула Тейлора vs Теорема Тейлора

Свойства степенных рядов (Теорема Абеля, радиус сходимости, ф-ма Коши-Адамара, равномерная сходимость внутри радиуса сходимости).

Теорема о голоморфности степенного ряда

Свойства голоморфных функций (голоморфность производной, гармоничность Re и Im, возможность восстановить одну из частей (Re или Im) с точностью до константы функций, пример на u(x, y) = x3 − 3xy 2 . Все без доказательств).

Теорема о характеристике полюса.

Теорема о связи коэффициентов в ряде Лорана с вычетами. Вычисление вычета.

Вычисление вещественных интегралов при помощи вычетов. (теорема о быстро убывающих интегралах, теорема Жордана)

\subsection{Дифференцируемость}

\begin{theorem}[О преобразовании касательного вектора]
Пусть $\gamma: [\alpha, \beta] \to D \subset \mathbb{C}$ --- гладкий путь, $t_0 \in (\alpha, \beta)$, и пусть функция $f: D \to \mathbb{C}$ голоморфна в точке $z_0 = \gamma(t_0)$. Тогда композиция $\Gamma = f \circ \gamma$ является гладким путём, и её производная в точке $t_0$ выражается формулой
\[
\Gamma'(t_0) = f'(z_0) \cdot \gamma'(t_0).
\]
Следовательно, голоморфное отображение $f$ действует на касательный вектор $\gamma'(t_0)$ как линейное преобразование плоскости, являющееся композицией поворота на угол $\arg f'(z_0)$ и растяжения в $|f'(z_0)|$ раз.
\end{theorem}

\begin{theorem}[О производной обратной функции]
Пусть $f\colon D_1\to D_2$, $D_1$, $D_2$ открыты, причём:
\begin{enumerate}
  \item $f$ --- биекция между $D_1$ и $D_2$
  \item $f\in\mathcal{O}(D_1)$ и $f'\neq 0$ на $D_1$
  \item $f^{-1}\in C(D_2)$
\end{enumerate}

 Тогда $f^{-1}\in\mathcal{O}(D_2)$, и для любой точки $w = f(z)$ верно:
\[
\bigl( f^{-1} \bigr)'(w) = \frac{1}{f'(z)} = \frac{1}{f'\bigl(f^{-1}(w)\bigr)}.
\]
\end{theorem}

\begin{theorem}[Условия Коши---Римана]
  Пусть $f\colon\mathbb{C}\supset D\to\mathbb{C},\ \Re f=u(x,y)$, $\Im f=v(x,y),\ z=x+iy,\ z_0=x_0+iy_0$ --- внутренняя точка для $D$.

  $\mathbb{C}$-дифференцируемость функции $f$ в $z_0$ равносильно двум условиям:
  \begin{enumerate}
    \item Функция $f$ $\mathbb{R}$-дифференцируемой в $z_0$.
    \item Выполнено \textbf{условие Коши---Римана}:
    \[ \begin{cases*}
      &u'_x(x_0,y_0)=v'_y(x_0,y_0)\\
      &u'_y(x_0,y_0)=-v'_x(x_0,y_0)\\
    \end{cases*} \]
  \end{enumerate}

  При этом $f'(z_0)=u'_x(x_0,y_0)+iv'_x(x_0,y_0)$.
\end{theorem}

\subsection{Коши}

\begin{lemma}[Гурса, или Коши-1]
Если функция $f$ голоморфна в области $D$, то её интеграл по границе любого треугольника $\Delta \subset D$ равен нулю:
\[
\int_{\partial \Delta} f(z) \, dz = 0.
\]
\end{lemma}

\begin{theorem}[Существование первообразной]
Пусть \( f\in\mathcal{O}(D) \), $D$ выпукло. Тогда в $D$ существует первообразная $f$.
\end{theorem}

\begin{theorem}[Оценка интеграла]
Пусть функция $f$ непрерывна вдоль кусочно-гладкого пути $\gamma: [\alpha, \beta] \to \mathbb{C}$. Тогда справедлива оценка
\[
\left| \int_{\gamma} f(z) \, dz \right| \leq \int_{\gamma} |f(z)| \, |dz|,
\]

В частности, если
\[
|f(z)| \leq M \quad \text{при всех} \quad z \in \gamma([\alpha, \beta]),
\]
то
\[
\left| \int_{\gamma} f(z) \, dz \right| \leq M \cdot |\gamma|,
\]
где $|\gamma|$ --- длина пути $\gamma$.
\end{theorem}

\begin{theorem}[Формула Ньютона–Лейбница]
Пусть $\gamma: [\alpha, \beta] \to D$ --- кусочно-гладкий путь в области $D$ и функция $f$ голоморфна в этой области. Обозначим через $\Phi$ первообразную $f$ вдоль $\gamma$. Тогда
\[
\int_{\gamma} f(z) \, dz = \Phi(\beta) - \Phi(\alpha).
\]
\end{theorem}

\begin{theorem}[Коши-2]
  Пусть \( f\in\mathcal{O}(D) \), $D$ выпукло:
  
  \begin{enumerate}
    \item Для замкнутого пути \( \gamma\subset D \):
    \[ \int_\gamma f\diff z=0 \]

    \item Для \( A,B\in D \) и \( \gamma_1,\ \gamma_2\subset D \) --- путей, соединяющих \( A,B \):
    \[ \int_{\gamma_1}f\diff z=\int_{\gamma_2}f\diff z \]
  \end{enumerate}
\end{theorem}

\begin{theorem}[Коши-3]
  Пусть \( f\in\mathcal{O}(D) \):
  
  \begin{enumerate}
    \item Пусть \( \gamma_1,\ \gamma_2\subset D \) --- замкнутые, не пересекающие друг друга пути, ориентированные положительно. Если часть плоскости между путями лежит в $D$:
    \[ \int_{\gamma_1}f\diff z=\int_{\gamma_2}f\diff z \]

    \item Пусть \( A,B\in D \) и \( \gamma_1,\ \gamma_2\subset D \) --- пути, соединяющие \( A,B \). Если часть плоскости между путями лежит в $D$:
    \[ \int_{\gamma_1}f\diff z=\int_{\gamma_2}f\diff z \]
  \end{enumerate}
\end{theorem}

\begin{theorem}[Интегральная формула Коши]
Пусть $D \subset \mathbb{C}$ --- область с кусочно-гладкой границей и функция $f$ голоморфна в некоторой области $G \supset \overline{D}$. Тогда для всех $z \in D$ справедлива формула
\[
f(z) = \frac{1}{2\pi i} \int_{\partial D} \frac{f(\zeta)}{\zeta - z} \, d\zeta.
\]
\end{theorem}

\subsection{Ряды Тейлора}

\begin{theorem}[Тейлора]
Пусть функция $f$ голоморфна в области $D \subset \mathbb{C}$ и $U_R(a) = \{ z \in \mathbb{C} : |z - a| < R \}$ — круг радиуса $R > 0$ с центром в точке $a \in D$, содержащийся в $D$. Введем обозначение
\[
c_n := \frac{1}{2\pi i} \int_{|\zeta - a| = r} \frac{f(\zeta) \, d\zeta}{(\zeta - a)^{n+1}}, \quad n = 0, 1, 2, \dots, \quad 0 < r < R.
\]
Числа $c_n$ не зависят от $r$ и называются \textbf{коэффициентами Тейлора} функции $f$ в точке $a$. Степенной ряд
\[
\sum_{n=0}^{\infty} c_n (z - a)^n
\]
называется \textbf{рядом Тейлора} функции $f$ с центром в точке $a$. Он сходится для всех $z \in U_R(a)$ и его сумма равна $f(z)$:
\[
f(z) = \sum_{n=0}^{\infty} c_n (z - a)^n \quad \text{при } |z - a| < R.
\]
\end{theorem}

\begin{theorem}[Неравенства Коши]
В условиях предыдущей теоремы при $0 < r < R$ и $n \in\mathbb{N}_0$ справедливы неравенства
\[
|c_n| \leq \frac{M(r)}{r^n}, \quad \text{где } M(r) := \max_{|\zeta-a|=r} |f(\zeta)|.
\]
\end{theorem}

\begin{theorem}[Лиувилля]
Пусть $f$ голоморфна во всей комплексной плоскости $\mathbb{C}$ и существует $M > 0$ такое, что
\[
|f(z)| \leq M
\]
для всех $z \in \mathbb{C}$.
Тогда $f(z) \equiv \text{const}$.
\end{theorem}

\subsection{Аналитичность}

\begin{theorem}[Свойства голоморфных функций]
Пусть $f$, $g\in\mathcal{O}(D)$. Тогда:

\begin{enumerate}
  \item Функция $f + g$ голоморфна в $D$, причём
  \[
  (f + g)'(z) = f'(z) + g'(z) \quad \text{для всех } z \in D.
  \]

  \item Функция $f \cdot g$ голоморфна в $D$, причём
  \[
  (f \cdot g)'(z) = f'(z) g(z) + f(z) g'(z) \quad \text{для всех } z \in D.
  \]

  \item Если $g(z) \neq 0$ для всех $z \in D$, то функция $f / g$ голоморфна в $D$, причём
  \[
  \left( \frac{f}{g} \right)'(z) = \frac{f'(z) g(z) - f(z) g'(z)}{[g(z)]^2} \quad \text{для всех } z \in D.
  \]
\end{enumerate}
\end{theorem}

\begin{theorem}[О бесконечной дифференцируемости голоморфной функции]
Функция $f$, голоморфная в произвольной области $D \subset \mathbb{C}$, имеет в $D$ производные всех порядков, которые также голоморфны в $D$.
\end{theorem}

\begin{lemma}[Вычисление коэффициентов ряда Тейлора]
Пусть функция $f$ голоморфна в окрестности точки $a \in \mathbb{C}$ и раскладывается в степенной ряд
\[
f(z) = \sum_{n=0}^{\infty} c_n (z - a)^n.
\]
Тогда коэффициенты $c_n$ выражаются через производные функции $f$ в точке $a$ по формуле
\[
c_n = \frac{f^{(n)}(a)}{n!}, \quad n \in\mathbb{N}_0
\]
\end{lemma}

\begin{theorem}[Морера]
Если функция $f$ непрерывна в области $D$ и интеграл от $f$ по границе любого треугольника $\Delta \subset D$ равен нулю, то $f$ голоморфна в $D$.
\end{theorem}

\begin{theorem}[Эквивалентные условия голоморфности]
Каждое из следующих условий эквивалентно голоморфности функции $f$ в точке $a \in \mathbb{C}$:

\begin{enumerate}
    \item Функция $f$ является $\mathbb{C}$-дифференцируемой в некоторой окрестности $U$ точки $a$;
    \item Функция $f$ аналитична в точке $a$, т.е. разлагается в степенной ряд с центром в точке $a$, сходящийся в некоторой окрестности $U$ точки $a$;
    \item Функция $f$ непрерывна в некоторой окрестности $U$ точки $a$ и интеграл от $f$ по границе любого треугольника $\Delta \subset U$ равен нулю.
\end{enumerate}
\end{theorem}

\subsection{Нули функции}

\begin{theorem}[Разложение голоморфной функции в нуле]
Пусть функция $f$ голоморфна в точке $a \in \mathbb{C}$ и $f(a) = 0$, но $f$ не равна тождественно нулю ни в какой окрестности точки $a$. Тогда в некоторой окрестности $U$ точки $a$ функция $f$ представляется в виде
\[
f(z) = (z - a)^n g(z),
\]
где $n$ — некоторое натуральное число, функция $g$ голоморфна в окрестности $U$ и отлична от нуля в $U$.
\end{theorem}

\begin{theorem}[Изолированность нулей голоморфной функции]
Пусть функция $f$ голоморфна в области $D$ и обращается в нуль в точке $a \in D$. Тогда верно одно из утверждений:
\begin{enumerate}
  \item $\exists U(a)\colon f \equiv 0$ в $U(a)$
  \item $\exists U(a)\colon f \neq 0$ в $\mathring{U}(a)$
\end{enumerate}
\end{theorem}

\begin{theorem}[Связь порядка нуля и разложения функции]
Пусть функция $f$ голоморфна в области $D$, $a \in D$ --- её изолированный ноль. Тогда следующие два утверждения эквивалентны:

\begin{enumerate}
    \item $\exists k \in \mathbb{N}\colon f(z) = (z-a)^k g(z)$, где $g$ --- голоморфная в $D$ функция, причём $g(a) \neq 0$.
    \item $f(a) = f'(a) = \dots = f^{(k-1)}(a) = 0$, но $f^{(k)}(a) \neq 0$.
\end{enumerate}
\end{theorem}

\subsection{Ряды Лорана}

\begin{theorem}[Лорана---Вейерштрасса]
Пусть функция $f(z)$ голоморфна в кольце
\[
V = \{ z \in \mathbb{C} : r < |z - a| < R \}
\]
с центром в точке $a \in \mathbb{C}$, где $0 \leq r < R \leq +\infty$. Положим
\[
c_n := \frac{1}{2\pi i} \int_{|\zeta - a| = \rho} \frac{f(\zeta) \, d\zeta}{(\zeta - a)^{n+1}} \quad \text{для всех } n \in \mathbb{Z} \text{ и } r < \rho < R.
\]
Числа $c_n$ не зависят от $\rho$ и называются \emph{коэффициентами Лорана} функции $f$ в кольце $V$. Ряд
\[
\sum_{n=-\infty}^{\infty} c_n (z - a)^n,
\]
называемый \emph{рядом Лорана} функции $f$ в кольце $V$, сходится для всех $z \in V$ и его сумма равна $f(z)$:
\[
f(z) = \sum_{n=-\infty}^{\infty} c_n (z - a)^n \quad \text{при } r < |z - a| < R.
\]
\end{theorem}

\subsection{Классификация особых точек}

\begin{theorem}[Критерий устранимой особой точки]
Для функции $f$, голоморфной в проколотой окрестности $V = \{0 < |z - a| < \varepsilon\}$ изолированной особой точки $a$, следующие утверждения эквивалентны:
\begin{enumerate}
    \item $\lim_{z\to a}f(z)\in\mathbb{C}$;
    \item коэффициенты главной части ряда Лорана равны нулю в некоторой $\mathring{U}(a)$;
    \item $f(z)$ ограничена в некоторой $\mathring{U}(a)$.
\end{enumerate}
\end{theorem}

\begin{theorem}[Критерий полюса]
Для функции $f$, голоморфной в проколотой окрестности $V = \{0 < |z - a| < \varepsilon\}$ изолированной особой точки $a$, следующие утверждения эквивалентны:
\begin{enumerate}
    \item $\lim_{z\to a}f(z)=\infty$;
    \item количество ненулевых коэффициентов в главной части ряда Лорана --- натуральное число;
    \item $\exists k\in\mathbb{N}\colon g(z)=(z-a)^kf(z)$ имеет $a$ устранимой особой точкой.
\end{enumerate}
\end{theorem}

\begin{theorem}[Критерий существенной особой точки]
Для функции $f$, голоморфной в проколотой окрестности $V = \{0 < |z - a| < \varepsilon\}$ изолированной особой точки $a$, следующие утверждения эквивалентны:
\begin{enumerate}
    \item $\nexists\lim_{z\to a}f(z)\in\overline{\mathbb{C}}$;
    \item количество ненулевых коэффициентов в главной части ряда Лорана бесконечно;
\end{enumerate}
\end{theorem}

\begin{theorem}[Сохоцкого]
Если $a \in \mathbb{C}$ --- существенно особая точка функции $f$, то для любого $A \in \overline{\mathbb{C}}$ можно найти последовательность точек $\{z_n\}$ такую, что $z_n \to a$ и
\[
\lim_{n \to \infty} f(z_n) = A.
\]
\end{theorem}

\begin{theorem}[Критерий бесконечности как особой точки]
  Точка $a = \infty$ является устранимой (полюсом, существенно особой) для функции $f(z)$ тогда и только тогда, когда точка $\zeta = 0$ является устранимой (полюсом, существенно особой) для функции $g(\zeta) := f(1/\zeta)$.

  Тогда $a = \infty$ есть:
  \begin{enumerate}
      \item \textbf{устранимая особая точка} функции $f$ $\iff$ $c_n = 0$ при всех $n \geq 1$;
      \item \textbf{полюс} функции $f$ $\iff$ существует $N \geq 1$ такое, что
      \[
      c_N \neq 0, \quad \text{но } c_n = 0 \quad \text{при } n \geq N + 1
      \]
      (число $N$ называется \emph{порядком полюса} в $\infty$);
      \item \textbf{существенно особая точка} функции $f$ $\iff$ $c_n \neq 0$ для бесконечного множества натуральных $n \geq 1$.
  \end{enumerate}
\end{theorem}

\subsection{Вычеты}

\begin{theorem}[Первая теорема о вычетах]
Пусть $D \subset \mathbb{C}$ --- область с кусочно-гладкой границей $\partial D$ и $G$ --- некоторая область в $\mathbb{C}$, содержащая замыкание $\overline{D}$ области $D$. Предположим, что функция $f$ голоморфна в области $G$, за исключением конечного числа изолированных особых точек $a_1, \ldots, a_n \in D$. Тогда
\[
\int_{\partial D} f(\zeta) \, d\zeta = 2\pi i \sum_{j=1}^{n} \operatorname{res}_{a_j} f.
\]
\end{theorem}

\begin{theorem}[Вторая теорема о вычетах]
Пусть функция $f$ голоморфна во всей расширенной комплексной плоскости $\overline{\mathbb{C}}$, за исключением конечного числа точек $\{a_\nu\}$. Тогда:
\[
\operatorname{res}_{\infty} f + \sum_{\nu} \operatorname{res}_{a_\nu} f = 0.
\]
\end{theorem}

\subsection{Принцип аргумента}

\begin{theorem}[О логарифмическом вычете]
Пусть выполняются следующие условия:
\begin{itemize}
    \item $\phi(z)$ --- голоморфная функция в области $D$,
    \item $\gamma \subset D$ --- замкнутый контур,
    \item $f(z)$ --- функция, голоморфная в $D$ и на контуре $\gamma$ (за исключением полюсов),
    \item Внутри контура $\gamma$ функция $f(z)$ имеет:
    \begin{itemize}
        \item Нули $z_1, \ldots, z_n$ порядков $\alpha_1, \ldots, \alpha_n$,
        \item Полюсы $\tilde{z}_1, \ldots, \tilde{z}_k$ порядков $\beta_1, \ldots, \beta_k$.
    \end{itemize}
\end{itemize}
Тогда справедливо равенство:
\[
\frac{1}{2\pi i} \int_{\gamma} \phi(z) \frac{f'(z)}{f(z)} \, dz = \sum_{i=1}^{n} \alpha_i \phi(z_i) - \sum_{j=1}^{k} \beta_j \phi(\tilde{z}_j).
\]
\end{theorem}

\begin{theorem}[Принцип аргумента]
Пусть функция $f(z)$ голоморфна в области $D$ (за исключением, быть может, конечного числа полюсов) и не имеет ни нулей, ни полюсов на замкнутом контуре $\gamma$.

Если взять общую теорему о логарифмическом вычете и подставить в неё $\phi(z) = 1$, то получим связь между интегралом и количеством нулей/полюсов:
\[
\frac{1}{2\pi i} \int_{\gamma} \frac{f'(z)}{f(z)} \, dz = N - P,
\]
где:
\begin{itemize}
    \item $N$ --- суммарное число нулей внутри контура $\gamma$ с учётом их порядков,
    \item $P$ --- суммарное число полюсов внутри контура $\gamma$ с учётом их порядков.
\end{itemize}

Так как интеграл слева равен $i \cdot \Delta_\gamma \arg f(z)$ (приращению аргумента вдоль контура), окончательная формула принципа аргумента выглядит так:
\[
\frac{1}{2\pi} \, \Delta_\gamma \arg f(z) = N - P.
\]
\end{theorem}

\begin{theorem}[Теорема Руше]
Пусть функции $f(z)$ и $\varphi(z)$ являются голоморфными в замкнутой области $\overline{G}$, причём на границе $\Gamma$ области $G$ имеет место неравенство
\[
|\varphi(z)| < |f(z)|, \quad z \in \Gamma.
\]
Тогда полное число нулей (с учётом кратностей) в области $G$ функции $F(z) = f(z) + \varphi(z)$ равно полному числу нулей функции $f(z)$.
\end{theorem}

\begin{theorem}[Основная теорема алгебры]
  Пусть дан полином над полем \( \mathbb{C} \):
  \[ P_n(z)=\sum_{i=0}^n a_iz^i,\quad n\geq 1,\ a_n\neq 0 \]
  
  Тогда \( P_n(z) \) имеет ровно \( n \) нулей с учётом кратности.
\end{theorem}